\chapter{Conclusão}\label{chap:conclusion}

\section{Objetivos Atingidos}\label{sec:goals}

\subsection{Objetivo Geral}

O método de simulação por QMDD foi implementado e integrado ao Ket, permitindo que o usuário escolha entre os diferentes métodos de simulação disponíveis na plataforma.

\subsection{Objetivos Específicos}

\begin{enumerate}[label=O\arabic*]
  \item Os algoritmos necessários para o QMDD foram compilados e descritos;

  \item Os algoritmos foram utilizados para implementar um simulador na linguagem de programação Rust;

  \item O novo simulador está integrado ao Ket, permitindo que usuários os utilizem de forma transparente e natural;

  \item O novo simulador foi testado em diferentes cenários, identificando sua vantagem sobre circuitos com poucas portas controladas.
\end{enumerate}

\section{Trabalhos Futuros}\label{sec:future}

% TODO: Rewrite this paragraph
A implementação atual utiliza vetores para armazenar os vértices criados e identificar repetições. Essa abordagem se apresentou suficiente para a aplicação correta dos algoritmos, porém apresenta complexidade linear em relação ao número de vértices já criados para criação de um novo vértice. Essa complexidade pode ser reduzida utilizando tabelas hash no lugar de vetores, reduzindo a complexidade para tempo constante. No momento, isso não é possível devido à forma como a implementação atual lida com números complexos, utilizando a implementação nativa de números decimais. Infelizmente, para que vértices possam ser identificados por tabelas hash, é necessária uma forma mais concisa de representar números irracionais, por exemplo o uso de números racionais extendidos pelos múltiplos de $\sqrt{2}$.

Além disso, a implementação atual está limitada ao comportamento do Ket de decompor portas arbitrárias em um conjunto fundamental de portas. Essa prática é desnecessária para o QMDD, que consegue naturalmente operar sobre matrizes arbitrárias, e impõe um custo adicional de tempo à simulação. Assim, uma melhoria futura seria a implementação de um simulador que não obrigue a decomposição de portas arbitrárias em portas fundamentais, promovendo ganho de desempenho.

% Quote LIMDD
Por fim, uma versão melhorada do QMDD foi proposta para conter o crescimento exponencial do número de vértices em circuitos com alta taxa de entrelaçamento. Essa versão, chamada LIMDD ({\itshape Local Invertible Map Decision Diagram}), utiliza a mesma ideia de compactação do QMDD, porém permite que cada aresta seja rotulada por uma matriz invertível localmente. Logo, é possível que os conhecimentos obtidos no estudo do QMDD sejam extendidos para o estudo do LIMDD, permitindo que o mesmo seja implementado e integrado ao Ket, promovendo ganho de desempenho em um conjunto mais amplo de circuitos quânticos. \cite{vinkhuijzen2023limdd}
